\documentclass{article}

\usepackage{threeparttable}
\usepackage{pgfplotstable}
\usepackage{booktabs}


\pgfplotstableset{row sep=newline,col sep=comma}

\providecommand{\emath}[1] {\ensuremath{#1}}
\newcounter {fncounter}
\setcounter {fncounter}{1}
\newcommand {\fnK}[1]      {\setcounter{fncounter}{#1}
                            \emath{\fnsymbol{fncounter}}}
\newcommand {\fnI}         {\fnK{1}}
\newcommand {\fnII}        {\fnK{2}}
\newcommand {\fnIII}       {\fnK{3}}

% Abbreviations
\newcommand{\with}         {w\textbackslash}               % Abbreviation w\
\newcommand{\gobblearg}[1]{}                               % Discard one argument

\newcommand{\includepgfplotstable}[2]{
\pgfplotstabletypeset[
  string type,
  header=false,
  skip first n=1,
  begin table={\begin{tabular}{lrrrrrr}
    \toprule
    #1
    \midrule\gobblearg},
  end table={
    \bottomrule
    \end{tabular}
  },
  every head row/.style={output empty row, after row={}},
  after row={\midrule},
  every last row/.style={after row={}}
]{#2}
}

\begin{document}

\begin{table*}[p]
\centering
\begin{threeparttable}
\centering
\begin{minipage}{\textwidth}

\includepgfplotstable{
        & &           & \# Cubes & \# Ambig 
          &           & \# Cubes \\
        & & \# Active & \with Ambig         & Grid 
          & \# Cubes & \with Config \\
        Dataset & Isovalue & cubes & Facets & Facets
            & \with MultiV & Changes \\
}{table_DC_multi_isov.csv}

\begin{tablenotes}
\item[\fnII] Isovalue for skin.
\item[\fnIII] Isovalue for skeleton.
\end{tablenotes}

\end{minipage}
\caption{Dual contouring multiple isosurface vertices per cube statistics.}
\end{threeparttable}
\end{table*}


\end{document}

